% Requires booktabs and graphicx. Supplementary material, not a replacement for
% original tables with different windows, references, or aggregation rules.
\begin{table}[t]
\centering\small
\caption{Constant-weight fusion controls on centered-square held-out overlap windows. Values are K24 terminal mean $E_u+E_p$ in percent. The train-fitted scalar minimizes the full-horizon squared field loss; the validation-selected scalar is a separate sensitivity control. Oracle weights use future reference trajectories and are diagnostic only.}
\label{tab:constant_alpha_supplement}
\begin{tabular}{lrr}
\toprule Method & S--H & H--P\\\midrule
Equal blend & 3.0321 & 1.6146\\
E2 probability blend & 2.1847 & 1.6618\\
Constant alpha (train fit) & 2.5590 & 1.1549\\
Constant alpha (validation sensitivity) & 1.7340 & 1.1607\\
T2-C (original checkpoint) & 1.1958 & 0.8465\\
Convex oracle (diagnostic) & 1.1842 & 0.8384\\
\bottomrule\end{tabular}
\end{table}

On the evaluated centered-square overlap windows, T2-C outperforms both
constant-weight controls at K24. This supports adaptive physical-space fusion
beyond the tested fixed mixing ratios. Attribution specifically to history
conditioning also requires the separate parameter-only input ablation.
The convex oracle uses future reference trajectories and is not deployable;
because it optimizes a full-horizon objective, it is not a strict lower bound
for every individual rollout step.

\begin{table}[t]
\centering\small
\caption{Circular Periodic supplement: K48 terminal errors on 12 aligned held-out windows. The common reference is the local POD-reconstructed field, not the full-order snapshot. Pressure is centered using area weights. Existing native auxiliary time/phase features are retained. Dense is a single-seed, total-parameter-matched control.}
\label{tab:dense_aligned_supplement}
\begin{tabular}{lrrr}
\toprule Method & $E_u$ (\%) & $E_p$ (\%) & $E_u+E_p$ (\%)\\\midrule
POD-Galerkin & 9.1217 & 36.2473 & 45.3691\\
Dense correction & 5.4590 & 20.1153 & 25.5743\\
Global MoE & 1.6072 & 5.8646 & 7.4719\\
RAL Periodic specialist & 1.9352 & 7.4275 & 9.3627\\
\bottomrule\end{tabular}
\end{table}

The structured local specialist improves K48 error over the tested
total-parameter-matched dense correction. However, Global MoE yields lower
terminal errors on these particular periodic windows. Thus, the results do
not establish uniform superiority of local specialists across regimes or
horizons. The global/local comparison also changes representation, mean-field
handling, and native auxiliary features, and should not be interpreted as a
single-factor causal test of locality.

\begin{figure}[t]
\centering
\includegraphics[width=\linewidth]{aligned_rollout_figures/four_method_error_growth.pdf}
\caption{Stepwise errors on the same 12 Circular Periodic held-out windows, with common POD-reconstructed reference and aligned pressure gauge. Logarithmic error axes are used. Step count is not a common physical duration across Reynolds numbers; per-Re native-time curves are supplied separately.}
\label{fig:aligned_error_growth_supplement}
\end{figure}

\begin{table}[t]
\centering\small
\caption{Native-implementation cost on one host (Xeon Platinum 8358P and RTX 4090): single-window K48 rollout, three warm-ups and ten synchronized repeats. The same first aligned window is used. Values exclude model/data loading, field decoding, and error computation. Neural models use GPU corrections; POD-Galerkin uses its native CPU NumPy path. GPU allocation is not total host memory. These are not full Top-2 pipeline timings.}
\label{tab:native_runtime_supplement}
\begin{tabular}{lrr}
\toprule Method & Median time (s) & Peak GPU allocated (MiB)\\\midrule
POD-Galerkin (CPU) & 0.3745 & 0.00\\
Dense correction & 0.5761 & 194.80\\
RAL Periodic specialist & 2.1264 & 194.78\\
Global MoE & 4.3295 & 194.78\\
\bottomrule\end{tabular}
\end{table}

Lower active parameter count does not imply lower latency: the dense control
is faster than the structured specialist in this native implementation,
despite evaluating more parameters. No general acceleration claim follows.
Our cross-regime claims concern evaluated regime transitions within
parameterized incompressible flows, not unseen physics or arbitrary regime
switching.

% Interpretation guide for existing ablations:
% DataOnly vs Full: physical backbone plus any native advancement changes.
% Vanilla-FNN-MoE vs Full: structured expert functions with routing retained.
% Global vs local: combined representation/model-decomposition effects.
% These controls do not individually prove that MoE itself causes all gains.
